\section{Art and Illusion in Meyrink}

\begin{quotex}
All that once was directly lived has become mere representation. \flright{\textsc{Guy Debord}, \textit{Society of the Spectacle}}

\end{quotex}
It is no secret that Plato harbored a distrust of the artist. There are two reasons for this. The first has to do state of mind of the artist, who creates an image of an image, an illusion purporting to be the form itself. \textbf{George Steiner}, in Grammars of Creation, poses these questions:

\begin{quotex}
How can the aesthetic represent, call to persuasive life that of which it possesses no direct, existential knowledge? A painter, wholly ignorant of seamanship, can depict a vessel performing some expert manoeuvre in raging seas. An arrant coward can sing famously of battle. A playwright who has never held public office will put in the mouths of his characters searching perceptions of statecraft. 

\end{quotex}
Beyond the illusion, there is the more serious ethical element. Steiner continues:

\begin{quotex}
Enactment in the arts, in fiction, are not only factitious and, in a fundamental sense, illusory. They are irresponsible. Devoid of authentic knowledge of that which it re-creates, the aesthetic plays with reality. The botched painting, the failed drama, are, in the strict sense of the word, inconsequential. The pilot who runs his craft onto the rocks, the loser in battle, are answerable even to the point of death. 

\end{quotex}
It is certainly banal to mention at this point the extent of this illusion of representation on contemporary Western life. Given Evola's fascination with \textbf{Gustav Meyrink}, I am more specifically interested in the mystical or spiritual literature, or in the many false currents claiming to represent Tradition today.

To start, we will look at a genuine work of spiritual value, such as Dante's Divine Comedy. As Dante himself points out, it is intended to be read on four levels: the \emph{literal}, \emph{allegorical}, \emph{moral}, and \emph{anagogical }meaning. The first three levels are clear enough and can be easily found; they are products of the work itself. Although the first three levels may evoke intense aesthetic feelings in the reader, whether of beauty or horror, that is an essentially passive process.

The anagogical, on the other hand, is of a different nature, as it must fully engage the reader. It is gnosis, pure and simple, which cannot simply be conveyed in words. The reader himself must reach the same levels of consciousness as Dante himself in order to grasp the meanings at this level. So far from being a work of illusion, an irresponsible re-creation, it really calls for a commitment of the reader, an appeal to make the same journey engaging the whole man. To follow that path is far from inconsequential.

Can we say the same of Meyrink? As Evola points out, Meyrink's works can be read simply as literature, although Meyrink himself claims they are a record of his own personal experiences. However, they don't have the gnostic, or anagogical, level as does the Divine Comedy. Rather, they are replete with visions, dreams, and trance states. Now Guenon made the rather mysterious claim\footnote{\url{https://www.gornahoor.net/?p=4559}} that he had to ``remediate certain evil consequences of \emph{The Green Face}," without providing specifics. More than likely, the consequences may be related to the confusion of the psychic and the spiritual as he describes in \emph{Reign of Quantity}. This confusion is based on the search for ``imagined powers" and ``extraordinary phenomena" instead of ``spiritual realization"; that is, there is the confusion between the ``forms of subtle manifestation" and the ``formless world".

It may not be so simple in fact. One of the characters in \emph{The Green Face} acknowledges the danger of hysteria associated with trances that leads downwards. He says, ``\emph{our} hysteria, on the other hand, is a matter of mental balance, the achieving of clarity, and is the way upwards, from insight through rational thought to knowledge through direct contemplation." This state is ``beyond error or even uncertainty".

There are other interesting notions such as ``to be master of one's own thoughts must also mean to be the all-powerful controller of one's own destiny" and how people ``take philosophy for a theory and not the practice." Meyrink writes about taming the body and taming one's thoughts. As for the former, it is insufficient, since he only makes a man a ``fakir". Regarding the latter, he writes:

\begin{quotex}
The next warriors your body will send out against you will be thoughts, whirring round you like a swarm of flies. Against them, the word of will-power is of no use. The more you strike out at them, the more furiously will they buzz around your head … It is impossible to command them to be still; there is but one way of escaping from them: by fleeing to a higher wakefulness. What you must do to achieve that, you must find out for yourself… The important thing is not to rid yourself of these thoughts for good, the purpose of your battle with them is to enter the state of higher wakefulness. 

\end{quotex}
Now we get to a crucial paragraph, which presumably provides the clue to understand all the visions and apparitions in Meyrink's works. These visions are merely appearances of something else. He continues:

\begin{quotex}
Once you have achieved that, the realm of ghosts, of which I have already spoken to you, will be near. You will see apparitions, both terrible and radiant, that will try to make you believe they are beings from another world. They are merely thoughts in visible form, over which you do not yet have power. The more sublime they appear, the more lethal they are, remember that!

Many a mistaken belief is founded on such apparitions, and has dragged mankind down into darkness. In spite of that, each one of these spectres conceals a deeper meaning. They are not mere pictures; irrespective of whether you can understand their symbolic language or not, they are signs of each stage of your spiritual development. 

\end{quotex}
So the psychic appearances that play such a prominent part in his books are not real in themselves, but illusions whose reality comes from the formless realm. Whether Meyrink has a ``direct, existential knowledge" of these teachings is hard to say. Regardless, there is not an anagogical level to the books nor an intiatic path that can be followed.



\flrightit{Posted on 2012-10-22 by Cologero }

\begin{center}* * *\end{center}

\begin{footnotesize}\begin{sffamily}



\texttt{Mihai on 2012-10-26 at 03:35 said: }

``Regardless, there is not an anagogical level to the books nor an intiatic path that can be followed."

I don't think that Meyrink ever intended these sort of things. As far as I can tell, he wrote about esoteric topics in an artistic way- just like a historian can present real historical events in a historical novel.

In any case, I have just begun reading The angel from the west window. The beginning is interesting. I will see what follows and then pass a judgement.


\hfill

\texttt{Cologero on 2012-10-26 at 12:31 said: }

I agree with you, Mihai, Meyrink did not consider his work to be an initiatic path. Rather it is a record of his own experiences. So it is a matter of how you judge those experiences.


\end{sffamily}\end{footnotesize}
